% Tabla D-T8. Generada automáticamente; no editar a mano.
% Requiere \usepackage{booktabs} (y graphicx si la tabla es ancha).
\begin{table}[htbp]
    \centering
    \caption{Submuestras utilizadas en los tres análisis.}
    \label{tab:D-T8}
    \small
    \resizebox{\textwidth}{!}{%
    \begin{tabular}{llp{0.40\textwidth}rr}
    \toprule
    Análisis & Unidad & Criterios de selección & $N_{\rm pl}$ & $N_\star$ \\
    \midrule
    I. Regímenes de radio & planeta & Descubiertos por tránsito, con radio, período, masa estelar y metalicidad & 4\,461 & 3\,390 \\
    I. Super-Tierras y sub-Neptunos & planeta & Tránsito, 1-4 $R_\oplus$ excluyendo 1.6-2.0 $R_\oplus$, con flujo incidente, masa estelar y metalicidad & 2\,565 & 1\,981 \\
    II. Propiedades físicas & planeta & 1-4 $R_\oplus$, masa verdadera y radio medido sin límites, flujo incidente, masa estelar y metalicidad & 574 & 385 \\
    III. Arquitectura & sistema & Al menos 2 planetas de 1-4 $R_\oplus$ y P $<$ 100 d, con masa estelar y metalicidad & 1\,513 & 615 \\
    \bottomrule
    \end{tabular}
    }
    \par\smallskip
    {\footnotesize\textit{Nota:} En el análisis III cada estrella es un sistema; el número de planetas es el total en esos sistemas.}
\end{table}
